\section{From the numerical invariant to a resolved family}
\label{sec:resolution}

We now assume that the central surface is not log canonical and that
the numerical argument has established constancy of
$\beta(X_t,0)=-P_{t,\mathrm{loc}}^2$.  Okuma's theorem gives a smooth
resolved family after a finite disk cover, with reduced nodal exceptional
restrictions.  We then use the fixed-ball Euler identity to show that
none of their nodes can smooth.  This yields the relative normal-crossing
divisor needed for complete transport in the next section.

A surface resolution is \emph{semigood} when its reduced exceptional
divisor has normal crossings; irreducible components may have self-nodes.
For a divisor in a smooth family of surfaces, \emph{relative normal
crossings} means that near its points there are coordinates $(z_1,z_2,u)$
in which the family map is $u$ and the divisor is $z_1=0$ or $z_1z_2=0$.

\begin{proposition}\label{prop:resolved-family}
Let $F(x,t)=f_t(x)$ satisfy the hypotheses of the main theorem.  Assume
that $X_0=\{f_0=0\}$ is not log canonical and that $\beta(X_t,0)$ is
constant.  After shrinking the parameter disk and making a finite base
change $t=u^e$, there is a proper modification
\[
 b:N\longrightarrow X'=\{F(x,u^e)=0\}
\]
such that $N\to\Delta_u$ is smooth, $b$ is an isomorphism off the origin
section, and its reduced exceptional divisor $E$ has relative normal
crossings.  The restriction $b_u$ is the minimal semigood resolution,
and $E|_{N_u}$ is its full reduced exceptional divisor.
\end{proposition}

\begin{proof}
The zero-fiber family is a flat deformation of a normal Gorenstein
surface singularity.  Each fiber has just the one singular point in the
chosen representative.  Thus constancy of the invariant at the section
is also constancy of the sum over the singularities of the fiber, as
used in Okuma's notation.

We first check the small-deformation setup in
\cite[Sections~3--4]{Okuma2004}.  Resolve the analytic total space, and
then principalize the pulled-back ideal of the section on the smooth
resolution \cite[Theorems~2.0.1 and~2.0.3, Lemma~4.0.3]{Wlodarczyk2009}.
The composite is proper and isomorphic off the section: the total space
is smooth there and the pulled-back ideal is the unit ideal there.
Its full reduced inverse image of the section is a simple normal crossing
divisor.  The exceptional support is proper over
the disk.  The critical loci for the parameter map on the resolution
and on all closed intersections of exceptional components lie in this
proper support.  These intersections are smooth.  On each of them a
critical component has constant parameter value; an intersection
contained in a fiber also has a single value.  Properness makes their
images closed analytic subsets of the disk.  After shrinking, their
only possible value is zero.  Hence the resolution over the punctured
disk is smooth, with relative normal crossing exceptional divisor.
These relative normal-crossing charts give a locally trivial deformation
of the exceptional divisor there.  This supplies the punctured-disk weak
simultaneous resolution required by Okuma's setup.

Okuma's Theorem~1.2 \cite{Okuma2004} now gives a simultaneous resolution
after a finite base change, with minimal semigood fiber resolutions and
a reduced divisor $E$ whose every restriction is the full reduced
exceptional divisor over the distinguished non-log-canonical section.
That section is our origin section, because there are no other
singularities.  Take the normalization of a branch of the finite base
change.  A local parameter then puts its map to the original disk in
the form $t=u^e$.  Flatness and smooth fibers make the resulting
resolved family smooth over this disk.  Formation of the Cartier
divisor restrictions commutes with this further base change.  In
particular their reducedness also holds at $u=0$.  The divisor $E$
contains no fiber, and its restrictions are effective Cartier
divisors; hence it is a relative effective Cartier divisor.  Its
support is exactly $b^{-1}(\{0\}\times\Delta_u)$.

It remains to prove relative normal crossings, since fiberwise normal
crossings alone do not give that conclusion.  Choose the fixed closed
Milnor ball $B$ from the preliminary Euler calculation and set
$N_B=b^{-1}(X'\cap(B\times\Delta_u))$.  Its projection to the disk
is a proper smooth submersion with boundary.  Near the sphere boundary
the modification is an isomorphism, and the required submersion follows
from the uniform boundary transversality of the original family.
Consequently $\chi((N_B)_u)$ is constant.  The complement of $E_u$ in
this fiber is isomorphic to $(X_{u^e}\cap B)\setminus\{0\}$.  The
preliminary identity $\chi(X_t\cap B)=1$ and additivity, using compact
supports on the complements, give
\[
 \chi(E_u)=\chi((N_B)_u),
\]
so the compact reduced nodal curves $E_u$ have constant Euler
characteristic.

Here is the local consequence of that constancy.  The divisor $E$ is
Cartier in the smooth total space $N$.  At a smooth point of $E_0$ it
is already smooth over the parameter.  At each node of $E_0$, the
holomorphic Morse lemma with parameter gives relative coordinates in
which its equation is
\[
 z_1z_2=a(u),\qquad a(0)=0.
\]
There are finitely many nodes.  Choose disjoint small neighborhoods of
them with transverse boundaries.  Smoothness over the parameter near
the remaining points of $E_0$, together with properness of $E\to\Delta_u$,
excludes new nodes outside these neighborhoods after shrinking the disk.
The remainder is a proper smooth family with boundary, so its topology
is unchanged.  Inside a node neighborhood a crossing has Euler
characteristic one, whereas a smoothing is an annulus and has Euler
characteristic zero.  Thus for general small $u$,
\[
 \chi(E_u)=\chi(E_0)-
 \#\{\text{nodes whose smoothing function }a\text{ is not identically zero}\}.
\]
No global labeling of components is needed; the calculation includes
self-intersection nodes.  Since the Euler characteristic is constant,
every smoothing function vanishes identically.  The resulting charts,
together with the charts at the smooth points, prove relative normal
crossings near $E_0$.  Properness of $E\to\Delta_u$ extends this
conclusion over a smaller disk.
\end{proof}

We have obtained the relative normal-crossing divisor needed for complete
transport.  The next section constructs the ambient homeomorphisms and
then removes the finite parameter cover.
